\documentclass{beamer}
\usepackage{xcolor}
\usepackage{graphicx, subcaption}
\usepackage{tikz}
\usepackage{booktabs}
\usepackage{listings}
\lstset{
    language=Python,
    basicstyle=\ttfamily\scriptsize,
    keywordstyle=\color{blue}\bfseries,
    commentstyle=\color{green!50!black},
    stringstyle=\color{red!70!black},
    showstringspaces=false,
    breaklines=true,
    frame=single,
    columns=fullflexible
}
\usetikzlibrary{trees,arrows.meta,positioning}
\usepackage{amsmath, amssymb, amsthm}

\title{MATH4907 Tutorial 2}
\author{Sanyang Liu}
\usetheme{Boadilla}
\useoutertheme{miniframes}
\setbeamertemplate{navigation symbols}{}

\newcommand{\bE}{\mathbb{E}}
\newcommand{\bQ}{\mathbb{Q}}
\newcommand{\bP}{\mathbb{P}}
\newcommand{\cF}{\mathcal{F}}
\newcommand{\R}{\mathbb{R}}
\newcommand{\order}{\mathcal{O}}
\newcommand{\cT}{\mathcal{T}}

\begin{document}
\AtBeginSection[]
{
    \begin{frame}{Table of Contents}
        \tableofcontents[
            currentsection,
            hideallsubsections
        ]
    \end{frame}
}
\maketitle

\begin{frame}{Personal Information}
\begin{itemize}
    \item TA: Sanyang Liu
    \item Office: Room 320A Run Run Shaw Building (Email me first)
    \item Email: u3013241@connect.hku.hk
\end{itemize}
\end{frame}

\begin{frame}{Table of Contents}
    \tableofcontents[]
\end{frame}

\section[A2Q1]{From Black--Scholes to Heat Equation}

\begin{frame}{Black--Scholes (BS) Equation}
Suppose we want to solve the BS equation for a European call option:
\[
\begin{cases}
V_t+\dfrac{1}{2}\sigma^2S^2V_{SS}
+rSV_S-rV=0,
& \forall (S,t)\in(0,\infty)\times[0,T),\\[0.25cm]
V(S,T)=(S-K)^+,
& \forall S\in(0,\infty),\\[0.25cm]
V(0,t)=0,
& \forall t\in[0,T),\\[0.25cm]
V(S,t)\sim S-Ke^{-r(T-t)},
& \text{as } S\to\infty,   \forall t\in[0,T).
\end{cases}
\]
\begin{block}{Numerical challenges}
\begin{enumerate}
    \item Backward in time.
    \item Coefficients depend on \(S\).
    \item Lower-order terms \(rSV_S\) and \(-rV\).
    \item Unbounded spatial domain.
\end{enumerate}
\end{block}
\end{frame}


\begin{frame}{Changes of Variables}
\footnotesize
\[
R(\tau_1)=\int_0^{\tau_1}r(s)\,ds,
\qquad
B(\tau_1)=\int_0^{\tau_1}
\left(r(s)-\frac12\sigma^2(s)\right)\,ds.
\]
\setlength{\tabcolsep}{2pt}
\renewcommand{\arraystretch}{1.0}
\begin{tabular}{p{0.27\textwidth}
                p{0.23\textwidth}
                p{0.45\textwidth}}
\toprule
\textbf{Motivation}
&
\textbf{Change of Variable}
&
\textbf{Equation}
\\
\midrule
Time reversal\\
backward $\to$ forward
&
$\tau_1=T-t$
&
$-V_{\tau_1}
+\frac12\sigma^2\textcolor{red}{S^2}V_{SS}
+rSV_S-rV=0$
\\[0.02cm]

Log transform\\
remove $S^2$
&
$y=\log S$
&
$-U_{\tau_1}
+\frac12\sigma^2U_{yy}
+\left(r-\frac12\sigma^2\right)U_y-r\textcolor{red}{U}=0$
\\[0.02cm]

Integrating factor\\
remove $U$
&
$v=e^{R(\tau_1)}U$
&
$-v_{\tau_1}
+\frac12\sigma^2v_{yy}
+\textcolor{red}{\left(r-\frac12\sigma^2\right)v_y}=0$
\\[0.02cm]

Moving coordinate\\
remove $v_y$
&
$x=y+B(\tau_1)$
&
$u_{\tau_1}=\textcolor{red}{\frac12\sigma^2}u_{xx}$
\\[0.02cm]

Time scaling\\
normalise diffusion
&
$\displaystyle
\tau=\frac12\int_0^{\tau_1}\sigma^2(s)\,ds$
&
$u_\tau=u_{xx}$
\\

\bottomrule
\end{tabular}
\end{frame}

\begin{frame}{Why Log Transform?}

The diffusion term in the Black--Scholes equation is $\frac{1}{2}\sigma^2S^2V_{SS}.$ To remove the \(S\)-dependence, assume a spatial transformation
\[
y=\phi(S),
\qquad
V(S,\tau)=U(y,\tau).
\]
By the chain rule,
\[
V_{SS}
=
U_{yy}\bigl(\phi'(S)\bigr)^2
+
U_y\phi''(S).
\]
Hence, the coefficient of \(U_{yy}\) becomes $\frac{1}{2}\sigma^2S^2\bigl(\phi'(S)\bigr)^2.$
We choose \(\phi\) such that $S^2\bigl(\phi'(S)\bigr)^2$ is independent of \(S\). Therefore,
\[
\phi'(S)\propto\frac{1}{S},
\qquad
\phi(S)=\log S
\]
up to an affine transformation. Thus we choose
\[
\boxed{y=\log S}.
\]
\end{frame}

\begin{frame}{Why Integrating Factor?}
After the time reversal and log transform, we obtain
\[
- U_{\tau_1}\textcolor{gray}{+
\frac{1}{2}\sigma^2U_{yy}
+
\left(r-\frac{1}{2}\sigma^2\right)U_y}
-rU=0.
\]
The zeroth-order term \(-rU\) is proportional to the unknown function.
To remove it without changing the spatial derivative terms, we try
\[
\phi(\tau_1)U(y,\tau_1)=v(y,\tau_1),
\]
where \(\phi\) is the integrating factor
\[
\boxed{
\phi(\tau_1)
=
\exp\left(\int_0^{\tau_1}r(s)\,ds\right)}.
\]
\end{frame}


\begin{frame}{Why Moving Coordinates?}

After removing the zeroth-order term, we obtain
\[
-v_{\tau_1}
\textcolor{gray}{+\frac{1}{2}\sigma^2v_{yy}
}+\left(r-\frac{1}{2}\sigma^2\right)v_y=0.
\]
For a transport equation, the solution moves along characteristics.
Therefore, introduce a moving coordinate
\[
x=y+B(\tau_1),
\qquad
B'=r-\frac{1}{2}\sigma^2.
\]
Write
\[
v(y,\tau_1)=u(x,\tau_1).
\]
Then
\[
v_{\tau_1}=u_{\tau_1}+B'(\tau_1)u_x,
\qquad
v_y=u_x,
\qquad
v_{yy}=u_{xx}.
\]
Therefore,
\[
u_{\tau_1}
=
\frac{1}{2}\sigma^2u_{xx}.
\]
\end{frame}




\section[A2Q2]{Finite Difference in 1-Dimensional}

\begin{frame}{1-Dimensional Stencil}
Suppose we want to approximate the \(m\)-th derivative \(f^{(m)}(x)\)
using the grid points
\[
x+\alpha_j h,\qquad j=1,\ldots,N.
\]
We seek coefficients \(c_j\) such that
\[
D_h f(x)
=
\frac{1}{h^m}
\sum_{j=1}^{N}
c_j f(x+\alpha_j h)
\]
approximates \(f^{(m)}(x)\) with a prescribed order. The point locations \(\{\alpha_j\}\) and coefficients \(\{c_j\}\)
together define \textcolor{blue}{a finite-difference stencil}.
\begin{block}{Questions}
How many points? Where should they be? Which coefficients?
What order of accuracy is possible?
\end{block}

\end{frame}

\begin{frame}{Taylor Expansion at Stencil Points}
Expand each value at the stencil point around \(x\):
\[
f(x+\alpha_j h)
=
\sum_{\ell\ge 0}
\frac{(\alpha_j h)^\ell}{\ell!}
f^{(\ell)}(x).
\]
Substituting these expansions into the definition of \(D_h f(x)\) gives
\[
D_h f(x)
=
\sum_{\ell\ge 0}
\frac{h^{\ell-m}}{\ell!}
\left(
\sum_{j=1}^{N}c_j\alpha_j^\ell
\right)
f^{(\ell)}(x).
\]
Define the moments of the stencil by
\[
M_\ell
:=
\sum_{j=1}^{N}c_j\alpha_j^\ell, \text{ then } D_h f(x)
=
\sum_{\ell\ge 0}
\frac{h^{\ell-m}}{\ell!}
M_\ell f^{(\ell)}(x).
\]

\end{frame}

\begin{frame}{Conditions for an Approximation of Order \(p\)}

To approximate \(f^{(m)}(x)\), the stencil coefficients must satisfy
\[
M_\ell=0,
\qquad \ell=0,\ldots,m-1,
\]
and
\[
M_m=m!.
\]
The first nonzero moment after \(M_m\) determines the order of accuracy.
If
\[
M_{m+1}=\cdots=M_{m+p-1}=0,
\qquad
M_{m+p}\neq 0,
\]
then
\[
D_hf(x)
=
f^{(m)}(x)
+
\frac{M_{m+p}}{(m+p)!}
h^p f^{(m+p)}(x)
+
O(h^{p+1}).
\]
Therefore, the finite-difference stencil has order of accuracy \(p\).

\end{frame}

\begin{frame}{The Coefficient-Matching Linear System}

For a chosen stencil, let
\[
\boldsymbol{c}
=
(c_1,\ldots,c_N)^T.
\]

The first \(m+1\) moment conditions give
\[
V\boldsymbol{c}
=
\boldsymbol{b},
\qquad
V\in\mathbb{R}^{(m+1)\times N},
\]
where
\[
V=
\begin{pmatrix}
1&\cdots&1\\
\alpha_1&\cdots&\alpha_N\\
\vdots&&\vdots\\
\alpha_1^m&\cdots&\alpha_N^m
\end{pmatrix},
\qquad
\boldsymbol{b}
=
\begin{pmatrix}
0\\
\vdots\\
0\\
m!
\end{pmatrix}.
\]
The matrix \(V\) is a generalized Vandermonde matrix.
\begin{itemize}
    \item The stencil locations determine \(V\).
    \item The derivative order determines \(\boldsymbol{b}\).
\end{itemize}
\end{frame}


\begin{frame}{Example (A2Q2): Setting the Stencil}

We want to approximate the second derivative \(f''(x)\).

Thus,
\[
m=2.
\]

Choose the symmetric three-point stencil
\[
\alpha=(-1,0,1),
\]
corresponding to
\[
x-h,\qquad x,\qquad x+h.
\]

Define
\[
D_hf(x)
=
\frac{1}{h^2}
\left[
c_{-1}f(x-h)
+c_0f(x)
+c_1f(x+h)
\right].
\]

The coefficients \(c_{-1},c_0,c_1\) are determined by
\[
M_0=0,\qquad M_1=0,\qquad M_2=2.
\]

\end{frame}

\begin{frame}{Example (A2Q2): Solving for the Coefficients}

The moment conditions give
\[
\begin{pmatrix}
1&1&1\\
-1&0&1\\
1&0&1
\end{pmatrix}
\begin{pmatrix}
c_{-1}\\c_0\\c_1
\end{pmatrix}
=
\begin{pmatrix}
0\\0\\2
\end{pmatrix}.
\]

Solving the linear system gives
\[
(c_{-1},c_0,c_1)=(1,-2,1).
\]

Therefore,
\[
D_hf(x)
=
\frac{
f(x-h)-2f(x)+f(x+h)
}{h^2}.
\]

\end{frame}

\begin{frame}{Example (A2Q2): Moments and Order of Accuracy}

For
\[
(c_{-1},c_0,c_1)=(1,-2,1),
\]
the relevant moments are
\[
M_0 =M_1=0, \,M_2=2,\, M_3=c_{-1}(-1)^3 + c_0 0^3 + c_1 1^3 = 0, \,M_4=2\neq0.
\]
Hence,
\[
D_hf(x)
=
\frac{M_2}{2!}f''(x)
+
\frac{M_4}{4!}h^2f^{(4)}(x)
+
O(h^4).
\]

Therefore,
\[
D_hf(x)
=
f''(x)
+\frac{1}{12}h^2f^{(4)}(x)
+O(h^4),
\]
and
\[
\boxed{
D_hf(x)=f''(x)+O(h^2).
}
\]
\end{frame}

\begin{frame}{A Systematic Procedure: 1D}

\small
\setlength{\itemsep}{2pt}

\begin{enumerate}
    \item Specify the target derivative order \(m\).

    \item Choose \(N\geq m+1\) distinct stencil locations
    \[
    (\alpha_1,\ldots,\alpha_N).
    \]

    \item Construct
    \[
    V\in\mathbb{R}^{N\times N},
    \qquad
    V_{ij}=\alpha_j^{\,i-1},
    \]
    and
    \[
    \boldsymbol b\in\mathbb{R}^N,
    \qquad
    b_i=m!\,\delta_{i,m+1}.
    \]

    \item Solve
    \[
    V\boldsymbol c=\boldsymbol b,
    \qquad
    \boldsymbol c\in\mathbb{R}^N.
    \]

    \item Compute
    \[
    M_\ell=\sum_{j=1}^N c_j\alpha_j^\ell,
    \qquad
    \ell=m+1,m+2,\ldots
    \]

    \item Find the first nonzero moment after \(M_m\), and determine
    the order of accuracy from the corresponding truncation error term.
\end{enumerate}
\end{frame}



\section[A2Q3]{Finite Difference in d-Dimensional}
\begin{frame}{Multi-Index Notation}
Let
\[
\boldsymbol{\alpha}
=
(\alpha_1,\ldots,\alpha_d)
\in\mathbb{N}_0^d
\]
denote the target derivative
\[
D^{\boldsymbol{\alpha}}f
=
\frac{\partial^{|\boldsymbol{\alpha}|}f}
{\partial x_1^{\alpha_1}\cdots\partial x_d^{\alpha_d}},
\]
where
\[
|\boldsymbol{\alpha}|
=
\alpha_1+\cdots+\alpha_d.
\]
For a multi-index \(\boldsymbol{\beta}=(\beta_1, \beta_2, \ldots, \beta_d)\), define
\[
\boldsymbol{\beta}!
=
\beta_1!\cdots\beta_d!,
\qquad
\boldsymbol{\xi}^{\boldsymbol{\beta}}
=
\prod_{i=1}^d \xi_i^{\beta_i}.
\]
\end{frame}

\begin{frame}{Multi-Dimensional Stencil}
Choose stencil points
\[
\boldsymbol{\xi}_j\in\mathbb{R}^d,
\qquad j=1,\ldots,N,
\]
and coefficients
\[
\boldsymbol{c}
=
(c_1,\ldots,c_N)^T.
\]

We define
\[
D_hf(\boldsymbol{x})
=
\frac{1}{h^{|\boldsymbol{\alpha}|}}
\sum_{j=1}^N
c_j f(\boldsymbol{x}+h\boldsymbol{\xi}_j).
\]

The point locations \(\boldsymbol{\xi}_j\) and coefficients \(c_j\)
define the multi-dimensional stencil.
\end{frame}


\begin{frame}{Multi-Index Taylor Expansion}
For each stencil point,
\[
f(\boldsymbol{x}+h\boldsymbol{\xi}_j)
=
\sum_{\boldsymbol{\beta}\in\mathbb{N}_0^d}
\frac{h^{|\boldsymbol{\beta}|}}
{\boldsymbol{\beta}!}
\boldsymbol{\xi}_j^{\boldsymbol{\beta}}
D^{\boldsymbol{\beta}}f(\boldsymbol{x}).
\]

Substituting into \(D_hf\) gives
\[
D_hf(\boldsymbol{x})
=
\sum_{\boldsymbol{\beta}}
\frac{h^{|\boldsymbol{\beta}|-|\boldsymbol{\alpha}|}}
{\boldsymbol{\beta}!}
M_{\boldsymbol{\beta}}
D^{\boldsymbol{\beta}}f(\boldsymbol{x}),
\]
where the multi-dimensional moments are
\[
M_{\boldsymbol{\beta}}
=
\sum_{j=1}^N
c_j\boldsymbol{\xi}_j^{\boldsymbol{\beta}}.
\]
\end{frame}

\begin{frame}{Moment Conditions in Multiple Dimensions}

To approximate \(D^{\boldsymbol{\alpha}}f\) with order \(p\), require
\[
M_{\boldsymbol{\beta}}
=
\begin{cases}
\boldsymbol{\alpha}!,
&
\boldsymbol{\beta}=\boldsymbol{\alpha},\\
0,
&
\boldsymbol{\beta}\neq\boldsymbol{\alpha},
\end{cases}
\]
for all multi-indices satisfying
\[
|\boldsymbol{\beta}|
\leq
|\boldsymbol{\alpha}|+p-1.
\]

The first nonzero moments after this range determine the
leading truncation error. Equivalently,
\[
D_hf(\boldsymbol{x})
=
D^{\boldsymbol{\alpha}}f(\boldsymbol{x})
+
O(h^p).
\]
\end{frame}

\begin{frame}{The Multivariate Vandermonde Matrix}

Choose multi-indices
\[
\boldsymbol{\beta}^{(1)},\ldots,\boldsymbol{\beta}^{(M)}
\]
and stencil points
\[
\boldsymbol{\xi}_1,\ldots,\boldsymbol{\xi}_N.
\]
The coefficient-matching system is
\[
V\boldsymbol{c}=\boldsymbol{b},
\]
where
\[
V=
\begin{pmatrix}
\boldsymbol{\xi}_1^{\boldsymbol{\beta}^{(1)}}&
\cdots&
\boldsymbol{\xi}_N^{\boldsymbol{\beta}^{(1)}}\\
\boldsymbol{\xi}_1^{\boldsymbol{\beta}^{(2)}}&
\cdots&
\boldsymbol{\xi}_N^{\boldsymbol{\beta}^{(2)}}\\
\vdots&&\vdots\\
\boldsymbol{\xi}_1^{\boldsymbol{\beta}^{(M)}}&
\cdots&
\boldsymbol{\xi}_N^{\boldsymbol{\beta}^{(M)}}
\end{pmatrix}
\in\mathbb{R}^{M\times N}.
\]
\end{frame}

\begin{frame}{Example (A2Q3): Approximate $f_{xy}$}
For A2Q3, the target derivative is
\[
D^{\boldsymbol{\alpha}}f=f_{xy},
\qquad
\boldsymbol{\alpha}=(1,1).
\]

Choose the four stencil points
\[
\boldsymbol{\xi}_1=(1,1),\quad
\boldsymbol{\xi}_2=(1,-1),\quad
\boldsymbol{\xi}_3=(-1,1),\quad
\boldsymbol{\xi}_4=(-1,-1).
\]

Use the multi-indices
\[
\boldsymbol{\beta}^{(1)}=(0,0),\quad
\boldsymbol{\beta}^{(2)}=(1,0),\quad
\boldsymbol{\beta}^{(3)}=(0,1),\quad
\boldsymbol{\beta}^{(4)}=(1,1).
\]

The matrix entries are defined by
\[
V_{ij}
=
\boldsymbol{\xi}_j^{\boldsymbol{\beta}^{(i)}}
=
(\xi_{j,1})^{\beta^{(i)}_1}
(\xi_{j,2})^{\beta^{(i)}_2}.
\]

For example,
\[
V_{2,1}
=
\boldsymbol{\xi}_1^{(1,0)}
=
1^1\,1^0=1.
\]
\end{frame}

\begin{frame}{Example (A2Q3): The Coefficient-Matching System}
Using $V_{ij}
=
\boldsymbol{\xi}_j^{\boldsymbol{\beta}^{(i)}}$, we obtain
\[
V=
\begin{pmatrix}
1&1&1&1\\
1&1&-1&-1\\
1&-1&1&-1\\
1&-1&-1&1
\end{pmatrix}.
\]

The moment conditions are
\[
V\boldsymbol{c}=\boldsymbol{b},
\qquad
b_i
=
\boldsymbol{\alpha}!\,
\delta_{\boldsymbol{\beta}^{(i)},\boldsymbol{\alpha}}.
\]
Thus,
\[
\begin{cases}
c_1+c_2+c_3+c_4=0,\\
c_1+c_2-c_3-c_4=0,\\
c_1-c_2+c_3-c_4=0,\\
c_1-c_2-c_3+c_4=1.
\end{cases}
\]
\end{frame}


\begin{frame}{Example (A2Q3): Coefficients and Accuracy}

Solving the linear system gives
\[
(c_1,c_2,c_3,c_4)
=
\left(
\frac14,-\frac14,-\frac14,\frac14
\right).
\]
Therefore,
\[
\begin{aligned}
&D_{h,k}f(x,y)\\
&=
\frac{
f(x+h,y+k)-f(x+h,y-k)
-f(x-h,y+k)+f(x-h,y-k)
}{4hk}.   
\end{aligned}
\]
The first neglected Taylor terms are of order
\[
O(h^2)+O(k^2).
\]
Hence,
\[
\boxed{
D_{h,k}f(x,y)
=
f_{xy}(x,y)+O(h^2+k^2).
}
\]
\end{frame}

\begin{frame}{A Systematic Procedure: 2D}

\small
\setlength{\itemsep}{2pt}

\begin{enumerate}
    \item Specify the target multi-index $\boldsymbol{\alpha}=(\alpha_1,\alpha_2).$
    \item Choose \(N\) stencil points
    \[
    \boldsymbol{\xi}_j
    =
    (\xi_{j,1},\xi_{j,2})
    \in\mathbb{R}^2,
    \qquad j=1,\ldots,N.
    \]

    \item Choose \(N\) multi-indices $\boldsymbol{\beta}^{(1)},\ldots,
    \boldsymbol{\beta}^{(N)}$ and construct
    \[
    V\in\mathbb{R}^{N\times N},
    \qquad
    V_{ij}
    =
    \boldsymbol{\xi}_j^{\boldsymbol{\beta}^{(i)}}
    =
    \xi_{j,1}^{\beta^{(i)}_1}
    \xi_{j,2}^{\beta^{(i)}_2}.
    \]
    and set $    b_i
    =
    \boldsymbol{\alpha}!\,
    \delta_{\boldsymbol{\beta}^{(i)},\boldsymbol{\alpha}},$
    \item Solve $V\boldsymbol c=\boldsymbol b,
    \qquad
    \boldsymbol c\in\mathbb{R}^N.$

\item Compute the multi-dimensional moments
\[
M_{\boldsymbol{\beta}}
=
\sum_{j=1}^N
c_j\boldsymbol{\xi}_j^{\boldsymbol{\beta}},
\qquad
|\boldsymbol{\beta}|\geq|\boldsymbol{\alpha}|,
\quad
\boldsymbol{\beta}\neq\boldsymbol{\alpha}.
\]
    \item Find the first nonzero moment after the target moment
    \(M_{\boldsymbol{\alpha}}\), and determine the order of accuracy.
\end{enumerate}
\end{frame}

\begin{frame}{A2Q3: Workflow}
Choose the four points $(x\pm h,y\pm k)$ and combine them with signs:
\[
\frac{f(x+h,y+k)-f(x+h,y-k)
-f(x-h,y+k)+f(x-h,y-k)}{4hk}.
\]
The sign pattern cancels the constant, first-order and pure second-order terms, while retaining $f_{xy}$.
\begin{block}{The message}
The mixed-derivative formula is not a separate trick. It is the same Taylor-plus-linear-system workflow in two variables.
\end{block}
\end{frame}


\section[A2Q4]{From Stencils to Numerical Schemes}

\begin{frame}{A Numerical Scheme Depends on the PDE}

Consider the heat equation
\[
u_\tau=u_{xx}.
\]

The PDE contains two differential operators:
\[
u_\tau
\qquad\text{and}\qquad
u_{xx}.
\]

To construct a numerical scheme, we need

\begin{itemize}
    \item a finite-difference stencil for \(u_\tau\);
    \item a finite-difference stencil for \(u_{xx}\);
    \item a rule for combining the two stencils.
\end{itemize}

\begin{block}{Key idea}
A numerical scheme is constructed from the differential operators appearing in the PDE.
\end{block}

\end{frame}


\begin{frame}{Finite-Difference Stencils}

For the time derivative,
\[
u_\tau(x_i,\tau_j)
=
\frac{u_i^{j+1}-u_i^j}{k}
+
O(k)
\qquad
\text{(forward difference)},
\]
and
\[
u_\tau(x_i,\tau_j)
=
\frac{u_i^j-u_i^{j-1}}{k}
+
O(k)
\qquad
\text{(backward difference)}.
\]

For the spatial second derivative,
\[
u_{xx}(x_i,\tau_j)
=
\frac{
u_{i-1}^j-2u_i^j+u_{i+1}^j
}{h^2}
+
O(h^2).
\]

\end{frame}


\begin{frame}{Explicit Scheme}

Using the forward difference in time and the central difference in space,
\[
u_\tau(x_i,\tau_j)
=
\frac{u_i^{j+1}-u_i^j}{k}
+
O(k),
\]
\[
u_{xx}(x_i,\tau_j)
=
\frac{
u_{i-1}^j-2u_i^j+u_{i+1}^j
}{h^2}
+
O(h^2).
\]

Substituting into \(u_\tau=u_{xx}\) gives
\[
\frac{u_i^{j+1}-u_i^j}{k}
=
\frac{
u_{i-1}^j-2u_i^j+u_{i+1}^j
}{h^2}
+
O(k)+O(h^2).
\]

Dropping the truncation-error terms gives the explicit scheme
\[
u_i^{j+1}
=
\lambda u_{i-1}^j
+
(1-2\lambda)u_i^j
+
\lambda u_{i+1}^j,
\qquad
\lambda=\frac{k}{h^2}.
\]

\end{frame}


\begin{frame}{Implicit Scheme}

Using the backward difference in time and the central difference
at the new time level,
\[
u_\tau(x_i,\tau_{j+1})
=
\frac{u_i^{j+1}-u_i^j}{k}
+
O(k),
\]
and
\[
u_{xx}(x_i,\tau_{j+1})
=
\frac{
u_{i-1}^{j+1}-2u_i^{j+1}+u_{i+1}^{j+1}
}{h^2}
+
O(h^2).
\]

The resulting scheme is
\[
\frac{u_i^{j+1}-u_i^j}{k}
=
\frac{
u_{i-1}^{j+1}-2u_i^{j+1}+u_{i+1}^{j+1}
}{h^2}.
\]

Equivalently,
\[
-\lambda u_{i-1}^{j+1}
+
(1+2\lambda)u_i^{j+1}
-\lambda u_{i+1}^{j+1}
=
u_i^j.
\]

\end{frame}


\begin{frame}{The \(\theta\)-Scheme}

The explicit and implicit schemes are special cases of
\[
\frac{u_i^{j+1}-u_i^j}{k}
=
(1-\theta)\delta_{xx}u_i^j
+
\theta\delta_{xx}u_i^{j+1},
\]
where
\[
\delta_{xx}u_i^j
=
\frac{
u_{i-1}^j-2u_i^j+u_{i+1}^j
}{h^2}.
\]

\begin{center}
\begin{tabular}{c|c}
\(\theta\) & Scheme \\ \hline
\(0\) & Explicit Euler \\
\(1\) & Implicit Euler \\
\(\frac12\) & Crank--Nicolson
\end{tabular}
\end{center}

\end{frame}


\begin{frame}{The Crank--Nicolson Scheme}

For \(\theta=\frac12\),
\[
\frac{u_i^{j+1}-u_i^j}{k}
=
\frac12\delta_{xx}u_i^j
+
\frac12\delta_{xx}u_i^{j+1}.
\]

Equivalently,
\[
-\frac{\lambda}{2}u_{i-1}^{j+1}
+
(1+\lambda)u_i^{j+1}
-\frac{\lambda}{2}u_{i+1}^{j+1}
\]
\[
=
\frac{\lambda}{2}u_{i-1}^{j}
+
(1-\lambda)u_i^{j}
+
\frac{\lambda}{2}u_{i+1}^{j},
\qquad
\lambda=\frac{k}{h^2}.
\]

The new time level appears on the left, so each time step
requires solving a linear system.

\end{frame}


\begin{frame}{Matrix Form}

Let
\[
U^j=(u_{1,j},\ldots,u_{n,j})^T
\]
denote the vector of interior values.

Define
\[
B=
\begin{pmatrix}
2&-1&&\\
-1&2&-1&\\
&\ddots&\ddots&\ddots\\
&&-1&2
\end{pmatrix}.
\]

Then the Crank--Nicolson scheme can be written as
\[
\left(I+\frac{\lambda}{2}B\right)U^{j+1}
=
\left(I-\frac{\lambda}{2}B\right)U^j
+
g^j,
\]
where \(g^j\) contains the boundary contributions.

\end{frame}


\begin{frame}{A2Q4: Definition of the Residual}

Let
\[
v_{i,j}=u(x_i,\tau_j)
\]
be the exact solution evaluated at the grid points.

For the Crank--Nicolson scheme, define
\[
\eta_{i,j}
=
v_{i,j+1}-v_{i,j}
-\frac{k}{2h^2}
\Big[
v_{i-1,j+1}-2v_{i,j+1}+v_{i+1,j+1}
\]
\[
\qquad\qquad\qquad
+
v_{i-1,j}-2v_{i,j}+v_{i+1,j}
\Big].
\]
The quantity \(\eta_{i,j}\) is the residual obtained by
substituting the exact solution into the discrete scheme.
\end{frame}


\begin{frame}{A2Q4: Taylor Expansion of the Residual}
Let
\[
x=x_i,\qquad \tau=\tau_j.
\]

The time expansion is
\[
u(x,\tau+k)-u(x,\tau)
=
ku_\tau+\frac{k^2}{2}u_{\tau\tau}
+O(k^3).
\]

For the spatial terms,
\[
u(x-h,\tau)-2u(x,\tau)+u(x+h,\tau)
=
h^2u_{xx}+O(h^4),
\]
and
\[
u(x-h,\tau+k)-2u(x,\tau+k)+u(x+h,\tau+k)
\]
\[
=
h^2u_{xx}
+
kh^2u_{xx\tau}
+
O(k^2h^2+h^4).
\]

Substituting these expansions into \(\eta_{i,j}\) gives
\[
\eta_{i,j}
=
k(u_\tau-u_{xx})
+
\frac{k^2}{2}
(u_{\tau\tau}-u_{xx\tau})
+
O(kh^2+k^3).
\]

\end{frame}

\begin{frame}{A2Q4: Local Truncation Error}

Since the exact solution satisfies the heat equation
\[
u_\tau=u_{xx},
\]
differentiating with respect to \(\tau\) gives
\[
u_{\tau\tau}=u_{xx\tau}.
\]

Therefore, the leading terms in the residual cancel:
\[
\eta_{i,j}
=
O(kh^2+k^3).
\]

The local truncation error is
\[
\cT_{i,j}
=
\frac{\eta_{i,j}}{k}.
\]

Hence,
\[
\cT_{i,j}
=
O(h^2+k^2).
\]

Therefore, the Crank--Nicolson scheme is
\[
\boxed{
\text{second-order accurate in space and time}.
}
\]

\end{frame}




\section[A2Q5]{Back to the Black--Scholes Equation}

\begin{frame}{The Transformed Black--Scholes Equation}

Let
\[
x=\log S,
\qquad
\tau=T-t.
\]

For constant \(r\) and \(\sigma\), the transformed equation is
\[
V_\tau
=
\frac{\sigma^2}{2}V_{xx}
+
\nu V_x
-rV,
\qquad
\nu=r-\frac{\sigma^2}{2}.
\]

For a European call option,
\[
V(x,0)=(e^x-K)^+.
\]

For the given parameters,
\[
S_0=K=100,
\qquad
r=0.05,
\qquad
\sigma=0.15,
\qquad
T=1,
\]
we have
\[
\nu=0.03875.
\]

\end{frame}


\begin{frame}{Explicit Finite-Difference Scheme}

At \((x_i,\tau_j)\), use
\[
V_\tau(x_i,\tau_j)
=
\frac{V_i^{j+1}-V_i^j}{k}
+
O(k),
\]
\[
V_x(x_i,\tau_j)
=
\frac{V_{i+1}^j-V_{i-1}^j}{2h}
+
O(h^2),
\]
and
\[
V_{xx}(x_i,\tau_j)
=
\frac{V_{i-1}^j-2V_i^j+V_{i+1}^j}{h^2}
+
O(h^2).
\]

Substituting into the transformed equation gives
\[
\frac{V_i^{j+1}-V_i^j}{k}
=
\frac{\sigma^2}{2}
\frac{V_{i-1}^j-2V_i^j+V_{i+1}^j}{h^2}
+
\nu
\frac{V_{i+1}^j-V_{i-1}^j}{2h}
-rV_i^j
+
O(k)+O(h^2).
\]

\end{frame}


\begin{frame}{Explicit Update Formula}

Dropping the truncation-error terms gives
\[
V_i^{j+1}
=
c_-V_{i-1}^j
+
c_0V_i^j
+
c_+V_{i+1}^j,
\]
where
\[
c_-
=
\frac{\sigma^2k}{2h^2}
-
\frac{\nu k}{2h},
\]
\[
c_0
=
1-\frac{\sigma^2k}{h^2}-rk,
\]
and
\[
c_+
=
\frac{\sigma^2k}{2h^2}
+
\frac{\nu k}{2h}.
\]

The initial values are
\[
V_i^0=(e^{x_i}-K)^+.
\]

\end{frame}


\begin{frame}{Boundary Conditions on a Truncated Domain}

The transformed spatial domain is unbounded:
\[
x\in\mathbb{R}.
\]

We truncate it to
\[
x\in[a,b].
\]

For a European call, use
\[
V(a,\tau)=0
\]
at the low-stock boundary, and
\[
V(b,\tau)=e^b-Ke^{-r\tau}
\]
at the high-stock boundary.

For \(a=-1\) and \(b=5\),
\[
V(-1,\tau)=0,
\]
\[
V(5,\tau)=e^5-100e^{-0.05\tau}.
\]

\end{frame}


\begin{frame}{Numerical Parameters}

Choose
\[
h=0.01,
\qquad
k=\frac{1}{225},
\qquad
m=225.
\]

The computation proceeds as follows:

\begin{enumerate}
    \item Construct the spatial grid \(x_i=a+ih\).
    \item Initialise
    \[
    V_i^0=(e^{x_i}-K)^+.
    \]
    \item Apply the explicit update for \(j=0,\ldots,m-1\).
    \item Impose the boundary values at every new time level.
    \item Evaluate the solution at \(x=\log S_0\).
\end{enumerate}

Since \(\log(100)\) is generally not a grid point,
use linear interpolation.

\end{frame}


\begin{frame}[fragile]{Implementation}

\begin{lstlisting}
import numpy as np

K = 100.0
S0 = 100.0
r = 0.05
sigma = 0.15

h = 0.01
m = 225
k = 1.0 / m
nu = r - sigma**2 / 2

c_minus = sigma**2*k/(2*h**2) - nu*k/(2*h)
c_zero  = 1 - sigma**2*k/h**2 - r*k
c_plus  = sigma**2*k/(2*h**2) + nu*k/(2*h)
\end{lstlisting}

\end{frame}

\begin{frame}[fragile]{Implementation (cont'd)}
\begin{lstlisting}
def value(a, b):
    n = round((b-a)/h)
    x = a + h*np.arange(n+1)

    V = np.maximum(np.exp(x)-K, 0.0)
    new_V = np.empty_like(V)

    for j in range(m):
        tau = (j+1)*k

        new_V[1:-1] = (
            c_minus*V[:-2]
            + c_zero*V[1:-1]
            + c_plus*V[2:]
        )
        new_V[0] = 0.0
        new_V[-1] = np.exp(b)-K*np.exp(-r*tau)
        V, new_V = new_V, V
    return np.interp(np.log(S0), x, V)
\end{lstlisting}
    
\end{frame}


\begin{frame}{Numerical Results}
On the truncated domain \([-1,5]\),
\[
V(\log 100,1)\approx 8.59888676.
\]
On the enlarged domain \([-2,10]\),
\[
V(\log 100,1)\approx 8.59888870.
\]

The difference is
\[
\left|
8.59888676-8.59888870
\right|
\approx
1.94\times10^{-6}.
\]

\begin{block}{Interpretation}
For these grid parameters, enlarging the computational domain has
a negligible effect on the computed option value.
\end{block}

\end{frame}




\section*{End}

\begin{frame}{}
\centering
\vfill
{\Huge\color{blue!70!black}\textbf{Thank You!}}
\vspace{1cm}
{\Large Questions?}
\vfill
\end{frame}

\end{document}
